\section{A common analytic interface and its PNT provider}\label{sec:analytic}

For $k\ge0$, put
\[
 \Pk_k=\{p\text{ prime}:2^k\le p<2^{k+1}\},
 \qquad
 A_k=\sum_{p\in\Pk_k}\frac1p.
\]
Fix once and for all a number
\begin{equation}\label{eq:mu-range}
 1<\mu<\log3.
\end{equation}

\begin{definition}[Common-threshold analytic supply]\label{def:AI}
The property $\mathbf{AI}(\mu)$ holds if there is an integer $K_{\mathrm{AI}}(\mu)$
such that, for every $k_0\ge K_{\mathrm{AI}}(\mu)$,
\begin{enumerate}[label=(\alph*)]
\item for every $k_0\le k\le3k_0$,
\begin{equation}\label{eq:AI-density}
 |\Pk_k|\ge\frac{2^k}{2\log(2^k)};
\end{equation}
\item the inclusive reciprocal window satisfies
\begin{equation}\label{eq:AI-mass}
 \sum_{k=k_0}^{3k_0}A_k\ge\mu.
\end{equation}
\end{enumerate}
The upper block index is inclusive, so
\begin{equation}\label{eq:inclusive-window}
 \bigcup_{k=k_0}^{3k_0}\Pk_k
 =\{p:2^{k_0}\le p<2^{3k_0+1}\}.
\end{equation}
\end{definition}

The common threshold and the endpoint in \eqref{eq:inclusive-window} are
load-bearing: density, reciprocal mass, and the later reservoir must be available
on the same block system.  The numerical choice $\mu=21/20$ is one convenient
admissible margin, but the construction uses only \eqref{eq:mu-range}.

Let $\pi(x)$ count primes at most $x$.  We use the prime number theorem in the
form
\begin{equation}\label{eq:PNT}
 \pi(x)\sim\frac{x}{\log x}\qquad(x\to\infty).
\end{equation}
A modern account is given by Soundararajan \cite{Soundararajan2006}; see also
Ingham, Theorem~23 \cite{Ingham1990}, and Hadamard's original proof
\cite{Hadamard1896}.

\begin{proposition}[Local reciprocal-prime law]\label{prop:local-law}
The prime number theorem implies
\begin{equation}\label{eq:local-x}
 \sum_{x\le p<2x}\frac1p
 =\frac{\log2}{\log x}
  +o\!\left(\frac1{\log x}\right).
\end{equation}
Consequently,
\begin{equation}\label{eq:local-k}
 A_k=\frac1k+o\!\left(\frac1k\right).
\end{equation}
\end{proposition}
\begin{proof}
Begin with $(x,2x]$.  Abel summation for the prime-counting step function and
$f(t)=1/t$ gives
\begin{equation}\label{eq:abel}
 \sum_{x<p\le2x}\frac1p
 =\frac{\pi(2x)}{2x}-\frac{\pi(x)}x
  +\int_x^{2x}\frac{\pi(t)}{t^2}\,dt.
\end{equation}
Set
\[
 \varepsilon(x)=\sup_{t\ge x}
 \left|\frac{\pi(t)\log t}{t}-1\right|.
\]
Then $\varepsilon(x)\to0$, and uniformly for $t\in[x,2x]$,
\[
 \pi(t)=\frac{t}{\log t}\bigl(1+O(\varepsilon(x))\bigr).
\]
Therefore
\begin{align*}
 \int_x^{2x}\frac{\pi(t)}{t^2}\,dt
 &=\int_x^{2x}\frac{dt}{t\log t}
   +O\!\left(\varepsilon(x)\int_x^{2x}\frac{dt}{t\log t}\right)\\
 &=\log\frac{\log(2x)}{\log x}
   +o\!\left(\frac1{\log x}\right)\\
 &=\frac{\log2}{\log x}
   +O\!\left(\frac1{\log^2x}\right)
   +o\!\left(\frac1{\log x}\right).
\end{align*}
The boundary term in \eqref{eq:abel} is
\[
 \frac1{\log(2x)}-\frac1{\log x}
 +O\!\left(\frac{\varepsilon(x)}{\log x}\right)
 =o\!\left(\frac1{\log x}\right).
\]
Changing $(x,2x]$ to $[x,2x)$ affects at most the endpoint terms, hence by
$O(1/x)=o(1/\log x)$.  At $x=2^k$, $k\ge2$, both endpoints are composite, so
the dyadic endpoint conventions agree exactly.  Since $\log(2^k)=k\log2$,
\eqref{eq:local-k} follows.
\end{proof}

\begin{proposition}[PNT supplies the common interface]\label{prop:PNT-AI}
For every fixed $\mu$ satisfying \eqref{eq:mu-range}, the local law implies
$\mathbf{AI}(\mu)$.
\end{proposition}
\begin{proof}
Choose constants
\[
 \frac1{2\log2}<c_D<1,
 \qquad
 \frac{\mu}{\log3}<c_M<1.
\]
Equation \eqref{eq:local-k} gives thresholds $K_D,K_M$ such that
$A_k\ge c_D/k$ for $k\ge K_D$ and $A_k\ge c_M/k$ for $k\ge K_M$.  Since
$1/p\le2^{-k}$ for $p\in\Pk_k$,
\[
 A_k\le\frac{|\Pk_k|}{2^k},
 \qquad
 |\Pk_k|\ge2^kA_k\ge\frac{2^k}{2k\log2},
\]
which is \eqref{eq:AI-density}.  Moreover,
\[
 \sum_{k=k_0}^{3k_0}\frac1k
 \ge\int_{k_0}^{3k_0+1}\frac{dt}{t}
 =\log\!\left(3+\frac1{k_0}\right)>\log3,
\]
so
\[
 \sum_{k=k_0}^{3k_0}A_k
 \ge c_M\sum_{k=k_0}^{3k_0}\frac1k
 >c_M\log3>\mu.
\]
Taking $K_{\mathrm{AI}}(\mu)=\max\{2,K_D,K_M\}$ proves both clauses at one
threshold.  This threshold is fixed before the construction chooses $k_0$.
\end{proof}

\begin{remark}
The Mertens asymptotic $\sum_{p\le x}1/p=\log\log x+B+o(1)$ is sufficient for a
fixed power window, but subtracting it across one dyadic block leaves
uncontrolled $o(1)$ errors while the local main term is only $O(1/\log x)$.
The PNT/Abel law supplies the required local scale directly; see Mertens
\cite{Mertens1874} for the historical asymptotic.
\end{remark}
